%%%%%%%%%%%%%%%%%%%%%%%%%%%%%%%%%%%%%%%%%%%%%%%%%%%%%%%%%%%%%%%
%                                                             %
% Cette version repart de 03_Draft\appendix\LE_definition.tex %
%                                                             %
%%%%%%%%%%%%%%%%%%%%%%%%%%%%%%%%%%%%%%%%%%%%%%%%%%%%%%%%%%%%%%%

\label{appsec:LE_def}
% More details on LE in France

In France, as in most OECD countries, separations and dismissals are classified into different types, depending on whether they are justified on economic or personal or disciplinary grounds. Given the scope of this paper, we focus on a panel of workers administratively registered as ``laid-off for economic reason'' (\textit{licenciement économique}, or LE) at the French employment agency, between 2002 and 2012. 

In France, dismissals for economic reasons require the employer to provide evidence of the existence of economic difficulties for the firm. Under French law, ``economic reasons'' entail four types of situations: (i) ``difficulties'' (persistent reduction of turnover, reduction of gross income, etc.), (ii) ``technological transformation'' (implying the need for different skills), (iii) ``re-organisation'' (in order to anticipate type-(i) issues, as opposed to improving performance) and (iv) ``cessation of business''. Economic dismissal can be individual or collective. During collective economic layoffs, the employer must also justify its choice of who remains employed and who is laid-off. 
French labor laws suggest employers to dismiss lower seniority workers first, as well as workers with lower family responsibility \citep{batut2020termination}. 

The process for ``licenciement économique'' depends on the size of the firm and the number of affected workers, with more stringent conditions regarding retraining imposed upon larger employer \citep{bender_margolis2002}. France being characterized by a high level of employee protection legislation (EPL)---in particular for permanent contracts \citep{oecd2004}---terminating a contract (as LE or another type of layoff) has economic consequences not only for the worker but also for the firm. For instance, laying-off workers for economic reasons imply severance payments.

% More details on UI system as required by R3 (or should we do a whole new appendix section?)

\paragraph{Unemployment insurance.} % Attention : info contemporaine. Any idea whether it was much different in our period of analysis?? 
Workers who lose their job involuntarily, whether being laid-off for economic reasons (LE), being displaced in plant closures, or being dismissed for reasons other than misconduct, are eligible for unemployment insurance benefits (Allocation d’aide au retour à l’emploi, ARE), provided they meet minimum contribution requirements and are actively searching for work. The unemployment insurance system does not differentiate between economic layoffs and other forms of involuntary separations in terms of benefit generosity or eligibility conditions. Benefits are proportional to previous earnings and typically replace between 57\% and 75\% of the worker’s reference wage, with an average replacement rate of around 70\% of prior net earnings. The duration of benefits depends on past employment history and age, and can extend up to approximately two years or more for older workers. In addition, unemployed workers are required to register with the public employment service (France Travail) and are subject to job-search requirements, with access to job-search assistance and training programs. Overall, while the French system provides substantial income insurance following job loss, it does not offer specific programs or more generous treatment targeted at workers laid off for economic reasons relative to other displaced workers. \textit{This implies that differences in post-displacement outcomes across groups are unlikely to be driven by differences in benefit generosity or access to labor market policies.}
